\section{Uncertainty paired designs and dependence}\label{sec:C26:uncertainty}
An agent evaluation has several sources of variation: task sampling, model sampling, environment nondeterminism, evaluator error, and adaptive state. Treating every model call as an independent observation ignores most of this structure. The appropriate uncertainty analysis follows the unit to which the claim generalizes. Repository-level generalization needs variation across repositories, not merely many attempts on the same few issues.

Paired comparisons reduce nuisance variation by evaluating systems on the same task instances. Define $\Delta_i=Y_{i,A}-Y_{i,B}$ and estimate its mean over the target population. Confidence intervals should preserve pairing. If tasks cluster by user, repository, or persistent session, resample clusters or use an explicitly justified hierarchical model. The choice should be declared before inspecting which interval is narrower.

Repeated seeds estimate stochastic execution variation conditional on the sampled tasks. They do not substitute for broader task coverage. Conversely, one run on many tasks can estimate a task-population average while leaving repeated-run reliability uncertain. Report both dimensions when the deployment depends on consistency. If compute prevents sufficient replication, narrow the claim rather than inventing stability from a positive mean.

Multiplicity matters when researchers compare many models, prompts, memory mechanisms, and metrics. Selecting the best configuration and reporting an ordinary confidence interval as if it had been fixed in advance understates uncertainty. Separate exploration from confirmation, report the selection procedure, and use untouched data for final comparison. A universal correction is not required for every descriptive analysis, but the inferential claim must reflect how the configuration was chosen.

Missing and censored outcomes need visible treatment. A timeout may indicate excessive latency, infrastructure failure, or an unfinished but promising trajectory. For a deadline-constrained task it is a failure to meet the deadline, regardless of possible eventual success. Additional analyses can examine completion-time distributions, but they should not erase the primary operational outcome. Report sensitivity to reasonable classification choices when the cause is ambiguous.

